\itemhead{p.11 练习 · 第3题（教材p.11）}{XJ-B1-C01-PR-P011-Y1526}
\label{XJ-B1-C01-PR-P011-Y1526}
{\footnotesize 湘教版教材p.11 / PDF第18页 / p.11 练习 / 第3题}
\field{原题}{已知集合 $A=\{(x,y)\mid x-y=0\}$，$B=\{(x,y)\mid x+y=0\}$，求 $A\cap B$。}
\field{参考答案}{解方程组得 $x=y=0$，故 $A\cap B=\{(0,0)\}$。}
{\small \field{关联知识}{S04 交集、并集及解集的联系（教材帮p.19、20 / PDF28、29）}
\field{方法与易错点}{先各自化简集合或解出约束，再在数轴上合并；有序数对交集通过联立方程求得。 易错：并集重复写元素；连接处一个闭端点也误删；把交集当并集。}
\field{答案依据}{答案册相关题组：PDF第4页（册内页3）；按当前题库原题逐项重算，不按旧版题号机械对套。}}
